\scanpage{27}
4. Supposons qu'on se donne une \unsure{isomor.|isom.}
\[
L\overset{\alpha}{\simeq}\underline{O}_S\qquad\text{i.e. une section }\inserted{\alpha^{-1}}\text{ de }\check L \tag*{(23)}
\]
alors la donnée d'un scindage équivaut
à la donnée d'une section de $t$ de $B$ au
dessus de $\alpha^{-1}$, La donnée d'un isom. (23)
et d'un scindage équivaut à la donnée d'une
trivialisation complète du fibré affine $E=E_\Omega$
i.e. d'un isom de fibrés affines
\[
E_\Omega\simeq E^1_S,
\]
ou encore à la donnée d'une section
$t$ de $B$ qui, en aucune \deleted{fibre} résiduelle,
n'est dans $\underline{O}_S\subset B$, \deleted{\unsure{Alors $b$ et $c$ s'identifient}}\inserted{i.e. telle que $(1,t)$}
soit une \uline{base} de $B$ sur $\underline{O}_S$. On peut \inserted{alors} identifier
$b$ et $c$ à des sections de $\underline{O}_S$, et celles-ci
sont définies par la condition (déduite de la
\unsure{dernière|deuxième} formule (13) en y faisant $u=v=t$)
\[
t^2-bt+c=0 \tag*{(24)}
\]
ou encore
\[
b=\operatorname{tr}t,\qquad c=N(t).\tag*{(25)}
\]
L'\deleted{formule} iso (23) définit \unsure{aussi}
\[
 L^{\otimes2}\overset{s^2}{\simeq}\underline{O}_S\qquad\deleted{\unsure{et}}
\]
et \deleted{\unsure{on y}} on aura alors
\[
\unsure{\mathscr D|\mathcal D|\mathfrak D}(\Omega/S)=(b^2-4c)s^2\qquad
\text{identifié, par la base }s^{\unsure{2}},\text{ à }b^2-4c \tag*{(26)}
\]

% archive sheet number: 26

